\section{功能性需求}

\subsection{执行者分析}

系统执行者按业务职责划分为用户、顾客、商家、骑手和管理员，用户表示使用平台公共能力的基础角色，顾客、商家、骑手和管理员是在用户基础上承担不同业务职责的具体角色，同一个自然人在通过相应审核后可以拥有多个业务角色，但每次业务操作仍按照当前角色及资源归属进行权限校验


\begin{longtable}{|L{0.10\linewidth}|L{0.15\linewidth}|L{0.74\linewidth}|}
\caption{系统执行者说明}\label{tab:actors}\\ \hline
\textbf{执行者} & \textbf{执行者关系} & \textbf{主要业务职责} \\ \hline
用户 & 基础执行者 & 注册和登录平台，维护个人信息，申请业务角色、查看通知 \\ \hline
顾客 & 特化角色 & 选购商品、管理订单与评价，使用智能点餐、个性化服务 \\ \hline
商家 & 特化角色 & 维护店铺及商品，处理订单，回复评价并查看经营数据 \\ \hline
骑手 & 特化角色 & 管理配送状态，接取并完成任务，处理配送过程中的异常 \\ \hline
管理员 & 特化角色 & 完成准入审核、平台治理、配送异常处置和平台数据查看 \\ \hline
\end{longtable}

\subsection{总用例}

\srsfigurelarge{figures/srs/09-system-use-case.png}{系统总用例图}{UML 用例图}{fig:system-use-case}
